\section{账本事件的区域制度深描：eastern-han}

\label{sec:ledger-depth-eastern-han}

本组叙事以本卷事件账本的可核年序为骨架，将政治节点放回地方资源、制度中介、人口流动与材料类型中。每条来源能够证明的范围不同，故不以朝廷命令、战报数字或后出传记替代基层社会的完整经验。

\subsection{地方资源与制度接口}

\eventanchor{EH-0025-GUANGWU-ACCESSION}{建武元年六月（25年）·刘秀在鄗即皇帝位，建元建武，重建汉朝名号}账本条目\texttt{EH-0025-GUANGWU-ACCESSION}记载，建武元年六月（25年），更始政权余部与刘秀集团发生“刘秀在鄗即皇帝位，建元建武，重建汉朝名号”。这一节点可放在更始政权在关中失序，河北军政集团需要以刘氏血统和建武年号整合征发与任官的条件下理解；其后续影响被账本概括为刘秀取得独立皇帝名义，战争目标由地方割据转为恢复汉廷并争夺全国承认。就地方尺度而言，事件并不只属于朝廷或统帅：征发、道路、仓储、户籍、市场、寺院与家庭关系都可能决定命令如何抵达、战事如何延长，或接收何以在不同地区呈现不一样的速度。

本条列举的核心材料为《后汉书·光武帝纪上》；《后汉纪》卷一；《资治通鉴》卷四十。这些材料较适合钉住事件的发生、官方表述或特定人物、地点，却不自动构成对全境人口、收入、兵数、信仰或动机的调查。账本的置信与材料提示是“即位年月可靠；即位仪式地点与参与者细节有异文”；来源限度进一步说明“光武建国、河北豪强与复汉合法性研究”。因此，不能把条目中的政权名称、行政分类或战报修辞直接转译为固定族群和统一社会意志，也不能将制度宣布写成各地同步实施。

在叙事上，更稳妥的做法是把这一事件视为一条需要地方中介才能运作的链：中央的命令、地方官与精英的选择、运输与劳力的条件、以及普通家庭可能面对的迁徙、赋役或安全风险彼此相连。现有材料能够显示其中若干环节，却保留了大量沉默。承认这种沉默并非放弃解释，而是避免以一条编年记录覆盖不同区域和社群的差异。

\eventanchor{EH-0025-LUOYANG-CAPITAL}{建武元年（25年）·刘秀定都洛阳，以雒阳为中央政府与军政调度中心}账本条目\texttt{EH-0025-LUOYANG-CAPITAL}记载，建武元年（25年），东汉发生“刘秀定都洛阳，以雒阳为中央政府与军政调度中心”。这一节点可放在河北根据地与关中战乱使长安难作稳定中心，洛阳更接近刘秀的军政支持网络的条件下理解；其后续影响被账本概括为洛阳成为东汉朝廷的制度与交通枢纽；其城市尺度和早期建设不能由后期遗址直接倒推。就地方尺度而言，事件并不只属于朝廷或统帅：征发、道路、仓储、户籍、市场、寺院与家庭关系都可能决定命令如何抵达、战事如何延长，或接收何以在不同地区呈现不一样的速度。

本条列举的核心材料为《后汉书·光武帝纪上》；《后汉纪》卷一；汉魏洛阳城遗址发掘报告。这些材料较适合钉住事件的发生、官方表述或特定人物、地点，却不自动构成对全境人口、收入、兵数、信仰或动机的调查。账本的置信与材料提示是“建都事实可靠；宫城复用和营建阶段须以考古分期核验”；来源限度进一步说明“东汉洛阳建都、都城复建与关东资源研究”。因此，不能把条目中的政权名称、行政分类或战报修辞直接转译为固定族群和统一社会意志，也不能将制度宣布写成各地同步实施。

在叙事上，更稳妥的做法是把这一事件视为一条需要地方中介才能运作的链：中央的命令、地方官与精英的选择、运输与劳力的条件、以及普通家庭可能面对的迁徙、赋役或安全风险彼此相连。现有材料能够显示其中若干环节，却保留了大量沉默。承认这种沉默并非放弃解释，而是避免以一条编年记录覆盖不同区域和社群的差异。

\eventanchor{EH-0026-GENGSHI-END}{建武二年（26年）·赤眉军入长安，更始帝刘玄败亡，关中更始政权终结}账本条目\texttt{EH-0026-GENGSHI-END}记载，建武二年（26年），赤眉军与更始政权发生“赤眉军入长安，更始帝刘玄败亡，关中更始政权终结”。这一节点可放在更始政权不能约束将领和地方豪强，关中供给、官爵分配与军事纪律持续恶化的条件下理解；其后续影响被账本概括为刘秀失去名义上的刘氏竞争者，但关中仍由赤眉和地方武装控制，统一战争继续。就地方尺度而言，事件并不只属于朝廷或统帅：征发、道路、仓储、户籍、市场、寺院与家庭关系都可能决定命令如何抵达、战事如何延长，或接收何以在不同地区呈现不一样的速度。

本条列举的核心材料为《后汉书·刘玄传》；《后汉纪》卷二；《资治通鉴》卷四十。这些材料较适合钉住事件的发生、官方表述或特定人物、地点，却不自动构成对全境人口、收入、兵数、信仰或动机的调查。账本的置信与材料提示是“核心过程可靠；刘玄死亡经过与赤眉内部决策有异文”；来源限度进一步说明“更始政权崩解、赤眉军与关中粮食危机研究”。因此，不能把条目中的政权名称、行政分类或战报修辞直接转译为固定族群和统一社会意志，也不能将制度宣布写成各地同步实施。

在叙事上，更稳妥的做法是把这一事件视为一条需要地方中介才能运作的链：中央的命令、地方官与精英的选择、运输与劳力的条件、以及普通家庭可能面对的迁徙、赋役或安全风险彼此相连。现有材料能够显示其中若干环节，却保留了大量沉默。承认这种沉默并非放弃解释，而是避免以一条编年记录覆盖不同区域和社群的差异。

\eventanchor{EH-0027-CHIMEI-DEFEAT}{建武三年（27年）·刘秀军在崤底一带受降赤眉，赤眉主力瓦解}账本条目\texttt{EH-0027-CHIMEI-DEFEAT}记载，建武三年（27年），东汉与赤眉军发生“刘秀军在崤底一带受降赤眉，赤眉主力瓦解”。这一节点可放在赤眉军在长安难以补给且内部失序，东汉军控制关东—关中通道后实施围堵与招降的条件下理解；其后续影响被账本概括为大规模流民军不再构成全国性竞争者，东汉可将兵力转向陇右、巴蜀和东部残余政权。就地方尺度而言，事件并不只属于朝廷或统帅：征发、道路、仓储、户籍、市场、寺院与家庭关系都可能决定命令如何抵达、战事如何延长，或接收何以在不同地区呈现不一样的速度。

本条列举的核心材料为《后汉书·光武帝纪上》；《后汉书·刘盆子传》；《后汉纪》卷二。这些材料较适合钉住事件的发生、官方表述或特定人物、地点，却不自动构成对全境人口、收入、兵数、信仰或动机的调查。账本的置信与材料提示是“年月大体可靠；战场位置和受降人数有争议”；来源限度进一步说明“赤眉覆灭、关中饥荒与流民武装研究”。因此，不能把条目中的政权名称、行政分类或战报修辞直接转译为固定族群和统一社会意志，也不能将制度宣布写成各地同步实施。

在叙事上，更稳妥的做法是把这一事件视为一条需要地方中介才能运作的链：中央的命令、地方官与精英的选择、运输与劳力的条件、以及普通家庭可能面对的迁徙、赋役或安全风险彼此相连。现有材料能够显示其中若干环节，却保留了大量沉默。承认这种沉默并非放弃解释，而是避免以一条编年记录覆盖不同区域和社群的差异。

\eventanchor{EH-0030-GONGSUN-SHU-EMPEROR}{建武六年（30年）·公孙述在成都称帝，国号成家，巴蜀形成独立政权中心}账本条目\texttt{EH-0030-GONGSUN-SHU-EMPEROR}记载，建武六年（30年），成家政权发生“公孙述在成都称帝，国号成家，巴蜀形成独立政权中心”。这一节点可放在巴蜀地形、粮赋与地方豪强网络使公孙述可拒绝关东政权的任命和征发的条件下理解；其后续影响被账本概括为东汉统一面临西南持久战，巴蜀的归属成为恢复帝国财政与交通的重要条件。就地方尺度而言，事件并不只属于朝廷或统帅：征发、道路、仓储、户籍、市场、寺院与家庭关系都可能决定命令如何抵达、战事如何延长，或接收何以在不同地区呈现不一样的速度。

本条列举的核心材料为《后汉书·公孙述传》；《华阳国志·公孙述志》；《后汉纪》卷五。这些材料较适合钉住事件的发生、官方表述或特定人物、地点，却不自动构成对全境人口、收入、兵数、信仰或动机的调查。账本的置信与材料提示是“称帝年份可靠；成家官制与实际控制范围有争议”；来源限度进一步说明“公孙述政权、巴蜀资源与区域国家形成研究”。因此，不能把条目中的政权名称、行政分类或战报修辞直接转译为固定族群和统一社会意志，也不能将制度宣布写成各地同步实施。

在叙事上，更稳妥的做法是把这一事件视为一条需要地方中介才能运作的链：中央的命令、地方官与精英的选择、运输与劳力的条件、以及普通家庭可能面对的迁徙、赋役或安全风险彼此相连。现有材料能够显示其中若干环节，却保留了大量沉默。承认这种沉默并非放弃解释，而是避免以一条编年记录覆盖不同区域和社群的差异。

\eventanchor{EH-0034-WEI-XIAO-SURRENDER}{建武十年（34年）·隗嚣死后，其子隗纯降汉，陇右主要割据力量被解除}账本条目\texttt{EH-0034-WEI-XIAO-SURRENDER}记载，建武十年（34年），东汉与陇右集团发生“隗嚣死后，其子隗纯降汉，陇右主要割据力量被解除”。这一节点可放在隗嚣在东汉、公孙述之间周旋，但陇右资源有限且其继承集团难以维持联盟的条件下理解；其后续影响被账本概括为东汉取得由关中通往河西、西南的战略通路，公孙述失去东部潜在盟友。就地方尺度而言，事件并不只属于朝廷或统帅：征发、道路、仓储、户籍、市场、寺院与家庭关系都可能决定命令如何抵达、战事如何延长，或接收何以在不同地区呈现不一样的速度。

本条列举的核心材料为《后汉书·隗嚣传》；《后汉纪》卷八；《资治通鉴》卷四十三。这些材料较适合钉住事件的发生、官方表述或特定人物、地点，却不自动构成对全境人口、收入、兵数、信仰或动机的调查。账本的置信与材料提示是“隗嚣卒年和隗纯降期可靠；降后处置细节有异文”；来源限度进一步说明“陇右割据、羌胡关系与汉军补给研究”。因此，不能把条目中的政权名称、行政分类或战报修辞直接转译为固定族群和统一社会意志，也不能将制度宣布写成各地同步实施。

在叙事上，更稳妥的做法是把这一事件视为一条需要地方中介才能运作的链：中央的命令、地方官与精英的选择、运输与劳力的条件、以及普通家庭可能面对的迁徙、赋役或安全风险彼此相连。现有材料能够显示其中若干环节，却保留了大量沉默。承认这种沉默并非放弃解释，而是避免以一条编年记录覆盖不同区域和社群的差异。

\eventanchor{EH-0035-EMANCIPATION-EDICT}{建武十一年（35年）·光武帝下诏限制将战乱中掠卖者继续为奴婢，并令核实释放}账本条目\texttt{EH-0035-EMANCIPATION-EDICT}记载，建武十一年（35年），东汉发生“光武帝下诏限制将战乱中掠卖者继续为奴婢，并令核实释放”。这一节点可放在长期战争造成掠卖、逃户和身份混乱，朝廷需要恢复编户、赋役和治安秩序的条件下理解；其后续影响被账本概括为国家以诏令重申自由民身份边界，但地方豪强和官府的执行不能由法令文字直接推定。就地方尺度而言，事件并不只属于朝廷或统帅：征发、道路、仓储、户籍、市场、寺院与家庭关系都可能决定命令如何抵达、战事如何延长，或接收何以在不同地区呈现不一样的速度。

本条列举的核心材料为《后汉书·光武帝纪下》；《后汉书·刘隆传》；《后汉纪》卷九。这些材料较适合钉住事件的发生、官方表述或特定人物、地点，却不自动构成对全境人口、收入、兵数、信仰或动机的调查。账本的置信与材料提示是“诏令存在可靠；实际释放规模和地方执行程度不明”；来源限度进一步说明“东汉初年奴婢、战乱掠卖与户籍恢复研究”。因此，不能把条目中的政权名称、行政分类或战报修辞直接转译为固定族群和统一社会意志，也不能将制度宣布写成各地同步实施。

在叙事上，更稳妥的做法是把这一事件视为一条需要地方中介才能运作的链：中央的命令、地方官与精英的选择、运输与劳力的条件、以及普通家庭可能面对的迁徙、赋役或安全风险彼此相连。现有材料能够显示其中若干环节，却保留了大量沉默。承认这种沉默并非放弃解释，而是避免以一条编年记录覆盖不同区域和社群的差异。

\eventanchor{EH-0036-CHENGJIA-FALL}{建武十二年（36年）·吴汉等攻入成都，公孙述战死，成家政权灭亡}账本条目\texttt{EH-0036-CHENGJIA-FALL}记载，建武十二年（36年），东汉与成家政权发生“吴汉等攻入成都，公孙述战死，成家政权灭亡”。这一节点可放在陇右平定后东汉得以集中兵力溯江入蜀，成家长期孤立且将领分歧加深的条件下理解；其后续影响被账本概括为巴蜀纳入东汉郡县任命与赋役网络，建武政权完成对主要竞争政权的军事压制。就地方尺度而言，事件并不只属于朝廷或统帅：征发、道路、仓储、户籍、市场、寺院与家庭关系都可能决定命令如何抵达、战事如何延长，或接收何以在不同地区呈现不一样的速度。

本条列举的核心材料为《后汉书·光武帝纪下》；《后汉书·公孙述传》；《华阳国志·公孙述志》。这些材料较适合钉住事件的发生、官方表述或特定人物、地点，却不自动构成对全境人口、收入、兵数、信仰或动机的调查。账本的置信与材料提示是“核心战事可靠；攻城过程与伤亡数字有争议”；来源限度进一步说明“东汉统一巴蜀、成都城防与战后地方整合研究”。因此，不能把条目中的政权名称、行政分类或战报修辞直接转译为固定族群和统一社会意志，也不能将制度宣布写成各地同步实施。

在叙事上，更稳妥的做法是把这一事件视为一条需要地方中介才能运作的链：中央的命令、地方官与精英的选择、运输与劳力的条件、以及普通家庭可能面对的迁徙、赋役或安全风险彼此相连。现有材料能够显示其中若干环节，却保留了大量沉默。承认这种沉默并非放弃解释，而是避免以一条编年记录覆盖不同区域和社群的差异。

\eventanchor{EH-0039-DUTIAN-SURVEY}{建武十五年（39年）·光武帝命郡国清查垦田、户口，因地方豪强隐匿与吏治争议引发纠纷}账本条目\texttt{EH-0039-DUTIAN-SURVEY}记载，建武十五年（39年），东汉发生“光武帝命郡国清查垦田、户口，因地方豪强隐匿与吏治争议引发纠纷”。这一节点可放在战后户籍散乱、田地兼并和地方官吏虚报削弱赋役基础，中央需重新掌握登记的条件下理解；其后续影响被账本概括为度田暴露中央财政目标与地方精英利益的冲突；诏令不等于获得可比的全国土地统计。就地方尺度而言，事件并不只属于朝廷或统帅：征发、道路、仓储、户籍、市场、寺院与家庭关系都可能决定命令如何抵达、战事如何延长，或接收何以在不同地区呈现不一样的速度。

本条列举的核心材料为《后汉书·光武帝纪下》；《后汉书·张湛传》；《后汉纪》卷十二。这些材料较适合钉住事件的发生、官方表述或特定人物、地点，却不自动构成对全境人口、收入、兵数、信仰或动机的调查。账本的置信与材料提示是“诏令与争议可靠；清查结果和全国覆盖率有争议”；来源限度进一步说明“东汉度田、户籍瞒报与豪强土地研究”。因此，不能把条目中的政权名称、行政分类或战报修辞直接转译为固定族群和统一社会意志，也不能将制度宣布写成各地同步实施。

在叙事上，更稳妥的做法是把这一事件视为一条需要地方中介才能运作的链：中央的命令、地方官与精英的选择、运输与劳力的条件、以及普通家庭可能面对的迁徙、赋役或安全风险彼此相连。现有材料能够显示其中若干环节，却保留了大量沉默。承认这种沉默并非放弃解释，而是避免以一条编年记录覆盖不同区域和社群的差异。

\eventanchor{EH-0040-TRUNG-REBELLION}{建武十六年（40年）·徵侧、徵贰在交趾起兵，攻取多处郡县，东汉岭南统治受挑战}账本条目\texttt{EH-0040-TRUNG-REBELLION}记载，建武十六年（40年），交趾徵氏集团与东汉发生“徵侧、徵贰在交趾起兵，攻取多处郡县，东汉岭南统治受挑战”。这一节点可放在交趾地方首领与太守苏定的治理冲突叠加，岭南的直接统治难以脱离地方网络维持的条件下理解；其后续影响被账本概括为东汉必须自岭南以外调兵平乱；越南后世记忆与汉文正史对行动者的表述须并列辨析。就地方尺度而言，事件并不只属于朝廷或统帅：征发、道路、仓储、户籍、市场、寺院与家庭关系都可能决定命令如何抵达、战事如何延长，或接收何以在不同地区呈现不一样的速度。

本条列举的核心材料为《后汉书·南蛮西南夷列传》；《后汉纪》卷十三；《大越史记全书》后出记载。这些材料较适合钉住事件的发生、官方表述或特定人物、地点，却不自动构成对全境人口、收入、兵数、信仰或动机的调查。账本的置信与材料提示是“起兵年份可靠；参与郡数和女性领导叙事的细节有争议”；来源限度进一步说明“徵氏起义、交趾地方首领与汉郡县研究”。因此，不能把条目中的政权名称、行政分类或战报修辞直接转译为固定族群和统一社会意志，也不能将制度宣布写成各地同步实施。

在叙事上，更稳妥的做法是把这一事件视为一条需要地方中介才能运作的链：中央的命令、地方官与精英的选择、运输与劳力的条件、以及普通家庭可能面对的迁徙、赋役或安全风险彼此相连。现有材料能够显示其中若干环节，却保留了大量沉默。承认这种沉默并非放弃解释，而是避免以一条编年记录覆盖不同区域和社群的差异。

\eventanchor{EH-0043-MA-YUAN-JIAOZHI}{建武十九年（43年）·马援平定交趾，徵氏集团覆灭，东汉重建岭南郡县控制}账本条目\texttt{EH-0043-MA-YUAN-JIAOZHI}记载，建武十九年（43年），东汉与交趾徵氏集团发生“马援平定交趾，徵氏集团覆灭，东汉重建岭南郡县控制”。这一节点可放在徵氏控制扩大使郡县官吏和交通节点失守，朝廷以水陆军和当地协力者恢复统治的条件下理解；其后续影响被账本概括为东汉恢复任命、征税与司法渠道，但后世铜柱传说不能单独证明当时边界或工程。就地方尺度而言，事件并不只属于朝廷或统帅：征发、道路、仓储、户籍、市场、寺院与家庭关系都可能决定命令如何抵达、战事如何延长，或接收何以在不同地区呈现不一样的速度。

本条列举的核心材料为《后汉书·马援传》；《后汉书·南蛮西南夷列传》；《后汉纪》卷十六。这些材料较适合钉住事件的发生、官方表述或特定人物、地点，却不自动构成对全境人口、收入、兵数、信仰或动机的调查。账本的置信与材料提示是“平乱年份可靠；作战路线、死伤和铜柱传说有争议”；来源限度进一步说明“马援南征、交趾行政与地方社会研究”。因此，不能把条目中的政权名称、行政分类或战报修辞直接转译为固定族群和统一社会意志，也不能将制度宣布写成各地同步实施。

在叙事上，更稳妥的做法是把这一事件视为一条需要地方中介才能运作的链：中央的命令、地方官与精英的选择、运输与劳力的条件、以及普通家庭可能面对的迁徙、赋役或安全风险彼此相连。现有材料能够显示其中若干环节，却保留了大量沉默。承认这种沉默并非放弃解释，而是避免以一条编年记录覆盖不同区域和社群的差异。

\subsection{道路、边地与人口移动}

\eventanchor{EH-0048-SOUTHERN-XIONGNU}{建武二十四年（48年）·南匈奴日逐王比归附汉廷，南北匈奴分裂，汉在边郡安置其部众}账本条目\texttt{EH-0048-SOUTHERN-XIONGNU}记载，建武二十四年（48年），东汉与南匈奴发生“南匈奴日逐王比归附汉廷，南北匈奴分裂，汉在边郡安置其部众”。这一节点可放在匈奴单于继承与资源分配冲突使比部寻求外援，东汉希望以附属部众缓冲北边的条件下理解；其后续影响被账本概括为南匈奴成为汉廷边防伙伴也保有自身首领网络，汉匈关系转为分裂政权间的长期协商。就地方尺度而言，事件并不只属于朝廷或统帅：征发、道路、仓储、户籍、市场、寺院与家庭关系都可能决定命令如何抵达、战事如何延长，或接收何以在不同地区呈现不一样的速度。

本条列举的核心材料为《后汉书·南匈奴传》；《后汉纪》卷二十一；《资治通鉴》卷四十四。这些材料较适合钉住事件的发生、官方表述或特定人物、地点，却不自动构成对全境人口、收入、兵数、信仰或动机的调查。账本的置信与材料提示是“归附年份可靠；部众数量、驻牧边界和自主程度有争议”；来源限度进一步说明“南匈奴内附、边郡安置与草原政治研究”。因此，不能把条目中的政权名称、行政分类或战报修辞直接转译为固定族群和统一社会意志，也不能将制度宣布写成各地同步实施。

在叙事上，更稳妥的做法是把这一事件视为一条需要地方中介才能运作的链：中央的命令、地方官与精英的选择、运输与劳力的条件、以及普通家庭可能面对的迁徙、赋役或安全风险彼此相连。现有材料能够显示其中若干环节，却保留了大量沉默。承认这种沉默并非放弃解释，而是避免以一条编年记录覆盖不同区域和社群的差异。

\eventanchor{EH-0057-MINGDI-ACCESSION}{建武中元二年（57年）·光武帝去世，太子刘庄即位，是为明帝}账本条目\texttt{EH-0057-MINGDI-ACCESSION}记载，建武中元二年（57年），东汉发生“光武帝去世，太子刘庄即位，是为明帝”。这一节点可放在建武时期主要竞争政权已被压制，刘庄作为长期确立的太子拥有相对稳定的继承基础的条件下理解；其后续影响被账本概括为东汉由开国军事动员转向守成朝廷，明帝朝的礼制、边疆和官僚政策在既有框架内调整。就地方尺度而言，事件并不只属于朝廷或统帅：征发、道路、仓储、户籍、市场、寺院与家庭关系都可能决定命令如何抵达、战事如何延长，或接收何以在不同地区呈现不一样的速度。

本条列举的核心材料为《后汉书·光武帝纪下》；《后汉书·明帝纪》；《后汉纪》卷二十八。这些材料较适合钉住事件的发生、官方表述或特定人物、地点，却不自动构成对全境人口、收入、兵数、信仰或动机的调查。账本的置信与材料提示是“继位年月可靠；光武晚年继承安排的细节有争议”；来源限度进一步说明“光武身后继承、宗室与东汉政治连续性研究”。因此，不能把条目中的政权名称、行政分类或战报修辞直接转译为固定族群和统一社会意志，也不能将制度宣布写成各地同步实施。

在叙事上，更稳妥的做法是把这一事件视为一条需要地方中介才能运作的链：中央的命令、地方官与精英的选择、运输与劳力的条件、以及普通家庭可能面对的迁徙、赋役或安全风险彼此相连。现有材料能够显示其中若干环节，却保留了大量沉默。承认这种沉默并非放弃解释，而是避免以一条编年记录覆盖不同区域和社群的差异。

\eventanchor{EH-0065-CHU-PRINCE-YING}{永平八年（65年）·楚王刘英因被告谋反获减死并被徙，其诏书涉及黄老与佛教供养的早期记录}账本条目\texttt{EH-0065-CHU-PRINCE-YING}记载，永平八年（65年），东汉发生“楚王刘英因被告谋反获减死并被徙，其诏书涉及黄老与佛教供养的早期记录”。这一节点可放在宗室王国仍拥有人员和资源，刘英案使宫廷以刑名与礼法限制其结交方术、僧侣和宾客的条件下理解；其后续影响被账本概括为该案提供早期汉地佛教书面线索，但不能由一纸诏令推断全国已有统一佛教制度。就地方尺度而言，事件并不只属于朝廷或统帅：征发、道路、仓储、户籍、市场、寺院与家庭关系都可能决定命令如何抵达、战事如何延长，或接收何以在不同地区呈现不一样的速度。

本条列举的核心材料为《后汉书·楚王英传》；《后汉书·明帝纪》；《牟子理惑论》相关传本。这些材料较适合钉住事件的发生、官方表述或特定人物、地点，却不自动构成对全境人口、收入、兵数、信仰或动机的调查。账本的置信与材料提示是“处分事实可靠；诏中佛教术语、文本层次与宗教实践范围有争议”；来源限度进一步说明“楚王英案、早期佛教语汇与宗室政治研究”。因此，不能把条目中的政权名称、行政分类或战报修辞直接转译为固定族群和统一社会意志，也不能将制度宣布写成各地同步实施。

在叙事上，更稳妥的做法是把这一事件视为一条需要地方中介才能运作的链：中央的命令、地方官与精英的选择、运输与劳力的条件、以及普通家庭可能面对的迁徙、赋役或安全风险彼此相连。现有材料能够显示其中若干环节，却保留了大量沉默。承认这种沉默并非放弃解释，而是避免以一条编年记录覆盖不同区域和社群的差异。

\eventanchor{EH-0073-DOU-GU-YIWULU}{永平十六年（73年）·窦固等出击北匈奴，取伊吾卢；班超奉使鄯善，东汉重启西域经营}账本条目\texttt{EH-0073-DOU-GU-YIWULU}记载，永平十六年（73年），东汉与北匈奴发生“窦固等出击北匈奴，取伊吾卢；班超奉使鄯善，东汉重启西域经营”。这一节点可放在北匈奴仍威胁河西交通，东汉试图以伊吾卢和绿洲盟友恢复通往西域的信息与补给路线的条件下理解；其后续影响被账本概括为汉廷重获进入西域的军事据点，但绿洲政权的服属需要持续外交和驻军，非一次远征所能固定。就地方尺度而言，事件并不只属于朝廷或统帅：征发、道路、仓储、户籍、市场、寺院与家庭关系都可能决定命令如何抵达、战事如何延长，或接收何以在不同地区呈现不一样的速度。

本条列举的核心材料为《后汉书·明帝纪》；《后汉书·窦固传》；《后汉书·班超传》。这些材料较适合钉住事件的发生、官方表述或特定人物、地点，却不自动构成对全境人口、收入、兵数、信仰或动机的调查。账本的置信与材料提示是“出师年份可靠；各路战果与班超使节行动的先后有争议”；来源限度进一步说明“永平北伐、伊吾卢与西域使团研究”。因此，不能把条目中的政权名称、行政分类或战报修辞直接转译为固定族群和统一社会意志，也不能将制度宣布写成各地同步实施。

在叙事上，更稳妥的做法是把这一事件视为一条需要地方中介才能运作的链：中央的命令、地方官与精英的选择、运输与劳力的条件、以及普通家庭可能面对的迁徙、赋役或安全风险彼此相连。现有材料能够显示其中若干环节，却保留了大量沉默。承认这种沉默并非放弃解释，而是避免以一条编年记录覆盖不同区域和社群的差异。

\eventanchor{EH-0075-CHEN-MU-CRISIS}{永平十八年（75年）·北匈奴攻车师，西域都护陈睦遇害，东汉一度撤出西域据点}账本条目\texttt{EH-0075-CHEN-MU-CRISIS}记载，永平十八年（75年），东汉、西域诸国与北匈奴发生“北匈奴攻车师，西域都护陈睦遇害，东汉一度撤出西域据点”。这一节点可放在东汉前推据点依赖河西补给和绿洲合作，北匈奴可利用车师地缘切断驻军联系的条件下理解；其后续影响被账本概括为东汉西域经营暂时收缩，显示都护名义与可持续的交通、粮秣和盟友并不等同。就地方尺度而言，事件并不只属于朝廷或统帅：征发、道路、仓储、户籍、市场、寺院与家庭关系都可能决定命令如何抵达、战事如何延长，或接收何以在不同地区呈现不一样的速度。

本条列举的核心材料为《后汉书·明帝纪》；《后汉书·耿恭传》；《后汉书·西域传》。这些材料较适合钉住事件的发生、官方表述或特定人物、地点，却不自动构成对全境人口、收入、兵数、信仰或动机的调查。账本的置信与材料提示是“核心危机可靠；撤军范围与北匈奴、车师各部责任有争议”；来源限度进一步说明“车师危机、都护制度与西域驻军补给研究”。因此，不能把条目中的政权名称、行政分类或战报修辞直接转译为固定族群和统一社会意志，也不能将制度宣布写成各地同步实施。

在叙事上，更稳妥的做法是把这一事件视为一条需要地方中介才能运作的链：中央的命令、地方官与精英的选择、运输与劳力的条件、以及普通家庭可能面对的迁徙、赋役或安全风险彼此相连。现有材料能够显示其中若干环节，却保留了大量沉默。承认这种沉默并非放弃解释，而是避免以一条编年记录覆盖不同区域和社群的差异。

\eventanchor{EH-0075-ZHANGDI-ACCESSION}{永平十八年（75年）·明帝去世，太子刘炟即位，是为章帝}账本条目\texttt{EH-0075-ZHANGDI-ACCESSION}记载，永平十八年（75年），东汉发生“明帝去世，太子刘炟即位，是为章帝”。这一节点可放在明帝在位时皇位继承已较明确，朝廷需要在边疆危机中保持中枢任命和财政连续的条件下理解；其后续影响被账本概括为章帝即位后由窦氏等外戚影响朝政，东汉宫廷权力结构进入新的调整期。就地方尺度而言，事件并不只属于朝廷或统帅：征发、道路、仓储、户籍、市场、寺院与家庭关系都可能决定命令如何抵达、战事如何延长，或接收何以在不同地区呈现不一样的速度。

本条列举的核心材料为《后汉书·明帝纪》；《后汉书·章帝纪》；《后汉纪》卷三十六。这些材料较适合钉住事件的发生、官方表述或特定人物、地点，却不自动构成对全境人口、收入、兵数、信仰或动机的调查。账本的置信与材料提示是“继位年月可靠；明帝遗命与辅政安排的细节有异文”；来源限度进一步说明“明章继承、外戚与守成政治研究”。因此，不能把条目中的政权名称、行政分类或战报修辞直接转译为固定族群和统一社会意志，也不能将制度宣布写成各地同步实施。

在叙事上，更稳妥的做法是把这一事件视为一条需要地方中介才能运作的链：中央的命令、地方官与精英的选择、运输与劳力的条件、以及普通家庭可能面对的迁徙、赋役或安全风险彼此相连。现有材料能够显示其中若干环节，却保留了大量沉默。承认这种沉默并非放弃解释，而是避免以一条编年记录覆盖不同区域和社群的差异。

\eventanchor{EH-0076-QIANG-REBELLION}{建初元年（76年）·西羌在陇西、金城等地反叛，东汉长期羌乱由此加剧}账本条目\texttt{EH-0076-QIANG-REBELLION}记载，建初元年（76年），东汉与西羌发生“西羌在陇西、金城等地反叛，东汉长期羌乱由此加剧”。这一节点可放在边郡官吏处置、徭役与牧地冲突叠加，羌部内部联盟又随战争重组的条件下理解；其后续影响被账本概括为陇右驻军、军费和迁徙压力持续增加；史书族名不应被视为固定而统一的社会实体。就地方尺度而言，事件并不只属于朝廷或统帅：征发、道路、仓储、户籍、市场、寺院与家庭关系都可能决定命令如何抵达、战事如何延长，或接收何以在不同地区呈现不一样的速度。

本条列举的核心材料为《后汉书·西羌传》；《后汉书·章帝纪》；《后汉纪》卷三十七。这些材料较适合钉住事件的发生、官方表述或特定人物、地点，却不自动构成对全境人口、收入、兵数、信仰或动机的调查。账本的置信与材料提示是“起事年份大体可靠；部落称谓、参与范围和初期战事有争议”；来源限度进一步说明“东汉羌乱、边郡征发与族群政治研究”。因此，不能把条目中的政权名称、行政分类或战报修辞直接转译为固定族群和统一社会意志，也不能将制度宣布写成各地同步实施。

在叙事上，更稳妥的做法是把这一事件视为一条需要地方中介才能运作的链：中央的命令、地方官与精英的选择、运输与劳力的条件、以及普通家庭可能面对的迁徙、赋役或安全风险彼此相连。现有材料能够显示其中若干环节，却保留了大量沉默。承认这种沉默并非放弃解释，而是避免以一条编年记录覆盖不同区域和社群的差异。

\eventanchor{EH-0079-WHITE-TIGER-HALL}{建初四年（79年）·章帝召集儒者在白虎观讨论五经异同，形成《白虎通》所保存的经义框架}账本条目\texttt{EH-0079-WHITE-TIGER-HALL}记载，建初四年（79年），东汉发生“章帝召集儒者在白虎观讨论五经异同，形成《白虎通》所保存的经义框架”。这一节点可放在官学中经师异说影响礼制和官僚教育，朝廷希望以皇帝主持的讨论稳定解释权的条件下理解；其后续影响被账本概括为经学语言更深地进入官方礼制与选官话语，但现存《白虎通》不能逐条等同会议原始记录。就地方尺度而言，事件并不只属于朝廷或统帅：征发、道路、仓储、户籍、市场、寺院与家庭关系都可能决定命令如何抵达、战事如何延长，或接收何以在不同地区呈现不一样的速度。

本条列举的核心材料为《后汉书·章帝纪》；《白虎通义》；《后汉纪》卷四十。这些材料较适合钉住事件的发生、官方表述或特定人物、地点，却不自动构成对全境人口、收入、兵数、信仰或动机的调查。账本的置信与材料提示是“会议存在可靠；《白虎通》与当时记录的编定关系有争议”；来源限度进一步说明“白虎观会议、经学国家化与文本编纂研究”。因此，不能把条目中的政权名称、行政分类或战报修辞直接转译为固定族群和统一社会意志，也不能将制度宣布写成各地同步实施。

在叙事上，更稳妥的做法是把这一事件视为一条需要地方中介才能运作的链：中央的命令、地方官与精英的选择、运输与劳力的条件、以及普通家庭可能面对的迁徙、赋役或安全风险彼此相连。现有材料能够显示其中若干环节，却保留了大量沉默。承认这种沉默并非放弃解释，而是避免以一条编年记录覆盖不同区域和社群的差异。

\eventanchor{EH-0088-HEDI-ACCESSION}{章和二年（88年）·章帝去世，皇子刘肇即位，是为和帝，窦太后临朝}账本条目\texttt{EH-0088-HEDI-ACCESSION}记载，章和二年（88年），东汉发生“章帝去世，皇子刘肇即位，是为和帝，窦太后临朝”。这一节点可放在新帝年少，窦氏凭皇太后身份和禁中网络取得辅政位置的条件下理解；其后续影响被账本概括为窦氏成为中枢人事与军事任命的核心，后续北伐功臣和外戚集团的关系更趋紧张。就地方尺度而言，事件并不只属于朝廷或统帅：征发、道路、仓储、户籍、市场、寺院与家庭关系都可能决定命令如何抵达、战事如何延长，或接收何以在不同地区呈现不一样的速度。

本条列举的核心材料为《后汉书·章帝纪》；《后汉书·和帝纪》；《后汉纪》卷四十九。这些材料较适合钉住事件的发生、官方表述或特定人物、地点，却不自动构成对全境人口、收入、兵数、信仰或动机的调查。账本的置信与材料提示是“继位事实可靠；章帝后期储位变动和窦氏权力配置有争议”；来源限度进一步说明“和帝继承、窦太后临朝与外戚政治研究”。因此，不能把条目中的政权名称、行政分类或战报修辞直接转译为固定族群和统一社会意志，也不能将制度宣布写成各地同步实施。

在叙事上，更稳妥的做法是把这一事件视为一条需要地方中介才能运作的链：中央的命令、地方官与精英的选择、运输与劳力的条件、以及普通家庭可能面对的迁徙、赋役或安全风险彼此相连。现有材料能够显示其中若干环节，却保留了大量沉默。承认这种沉默并非放弃解释，而是避免以一条编年记录覆盖不同区域和社群的差异。

\eventanchor{EH-0089-DOU-XIAN-JILU}{永元元年（89年）·窦宪等北出稽落山，大破北匈奴，立燕然山铭纪功}账本条目\texttt{EH-0089-DOU-XIAN-JILU}记载，永元元年（89年），东汉与北匈奴发生“窦宪等北出稽落山，大破北匈奴，立燕然山铭纪功”。这一节点可放在窦氏希望以大规模北伐巩固辅政威望，南匈奴与边郡骑兵为行动提供协力的条件下理解；其后续影响被账本概括为北匈奴受到重创，窦宪军功显著抬升；铭文是胜利宣示，不能直接作为战果总数的独立证明。就地方尺度而言，事件并不只属于朝廷或统帅：征发、道路、仓储、户籍、市场、寺院与家庭关系都可能决定命令如何抵达、战事如何延长，或接收何以在不同地区呈现不一样的速度。

本条列举的核心材料为《后汉书·和帝纪》；《后汉书·窦宪传》；《后汉书·南匈奴传》；燕然山铭。这些材料较适合钉住事件的发生、官方表述或特定人物、地点，却不自动构成对全境人口、收入、兵数、信仰或动机的调查。账本的置信与材料提示是“战役年月可靠；兵数、斩获与铭文功绩措辞有夸饰可能”；来源限度进一步说明“窦宪北伐、燕然山铭与东汉边疆纪功研究”。因此，不能把条目中的政权名称、行政分类或战报修辞直接转译为固定族群和统一社会意志，也不能将制度宣布写成各地同步实施。

在叙事上，更稳妥的做法是把这一事件视为一条需要地方中介才能运作的链：中央的命令、地方官与精英的选择、运输与劳力的条件、以及普通家庭可能面对的迁徙、赋役或安全风险彼此相连。现有材料能够显示其中若干环节，却保留了大量沉默。承认这种沉默并非放弃解释，而是避免以一条编年记录覆盖不同区域和社群的差异。

\eventanchor{EH-0091-BAN-CHAO-PROTECTOR}{永元三年（91年）·班超受命为西域都护，东汉在多处绿洲重新建立使节与军事网络}账本条目\texttt{EH-0091-BAN-CHAO-PROTECTOR}记载，永元三年（91年），东汉与西域诸国发生“班超受命为西域都护，东汉在多处绿洲重新建立使节与军事网络”。这一节点可放在北匈奴势力受挫后，东汉试图将临时使团关系转化为较稳定的都护协调的条件下理解；其后续影响被账本概括为汉廷恢复对部分绿洲的影响，但该体系依赖地方王权、商路和补给，不能表述为完整的直接统治。就地方尺度而言，事件并不只属于朝廷或统帅：征发、道路、仓储、户籍、市场、寺院与家庭关系都可能决定命令如何抵达、战事如何延长，或接收何以在不同地区呈现不一样的速度。

本条列举的核心材料为《后汉书·班超传》；《后汉书·西域传》；《后汉纪》卷五十二。这些材料较适合钉住事件的发生、官方表述或特定人物、地点，却不自动构成对全境人口、收入、兵数、信仰或动机的调查。账本的置信与材料提示是“任命年份可靠；各国服属程度和驻军规模有争议”；来源限度进一步说明“班超都护、西域绿洲政治与汉军驻防研究”。因此，不能把条目中的政权名称、行政分类或战报修辞直接转译为固定族群和统一社会意志，也不能将制度宣布写成各地同步实施。

在叙事上，更稳妥的做法是把这一事件视为一条需要地方中介才能运作的链：中央的命令、地方官与精英的选择、运输与劳力的条件、以及普通家庭可能面对的迁徙、赋役或安全风险彼此相连。现有材料能够显示其中若干环节，却保留了大量沉默。承认这种沉默并非放弃解释，而是避免以一条编年记录覆盖不同区域和社群的差异。

\subsection{社会关系与行政分类}

\eventanchor{EH-0092-DOU-CLAN-PURGE}{永元四年（92年）·和帝联同宦官郑众诛除窦宪集团，结束窦太后家族的军政主导}账本条目\texttt{EH-0092-DOU-CLAN-PURGE}记载，永元四年（92年），东汉发生“和帝联同宦官郑众诛除窦宪集团，结束窦太后家族的军政主导”。这一节点可放在窦宪以北伐和禁军权势压过皇帝，和帝需借近侍和宫廷通道打破外戚控制的条件下理解；其后续影响被账本概括为和帝得以亲政，宦官因参与政变进入更关键的政治位置，外戚—宦官竞争的结构加深。就地方尺度而言，事件并不只属于朝廷或统帅：征发、道路、仓储、户籍、市场、寺院与家庭关系都可能决定命令如何抵达、战事如何延长，或接收何以在不同地区呈现不一样的速度。

本条列举的核心材料为《后汉书·和帝纪》；《后汉书·窦宪传》；《后汉书·宦者列传》。这些材料较适合钉住事件的发生、官方表述或特定人物、地点，却不自动构成对全境人口、收入、兵数、信仰或动机的调查。账本的置信与材料提示是“政变事实可靠；和帝自主程度与郑众角色有争议”；来源限度进一步说明“和帝亲政、窦氏覆灭与宦官近侍研究”。因此，不能把条目中的政权名称、行政分类或战报修辞直接转译为固定族群和统一社会意志，也不能将制度宣布写成各地同步实施。

在叙事上，更稳妥的做法是把这一事件视为一条需要地方中介才能运作的链：中央的命令、地方官与精英的选择、运输与劳力的条件、以及普通家庭可能面对的迁徙、赋役或安全风险彼此相连。现有材料能够显示其中若干环节，却保留了大量沉默。承认这种沉默并非放弃解释，而是避免以一条编年记录覆盖不同区域和社群的差异。

\eventanchor{EH-0097-GAN-YING-ANXI}{永元九年（97年）·班超遣甘英西使，至安息属地而未能到达大秦}账本条目\texttt{EH-0097-GAN-YING-ANXI}记载，永元九年（97年），东汉、安息与西域诸国发生“班超遣甘英西使，至安息属地而未能到达大秦”。这一节点可放在东汉西域网络暂获恢复，班超希望获得更远西方政权和海陆路线的情报的条件下理解；其后续影响被账本概括为汉文史料保存了经安息获得的西方地理知识，但旅行未达终点，不能据此推定汉罗直接外交关系。就地方尺度而言，事件并不只属于朝廷或统帅：征发、道路、仓储、户籍、市场、寺院与家庭关系都可能决定命令如何抵达、战事如何延长，或接收何以在不同地区呈现不一样的速度。

本条列举的核心材料为《后汉书·西域传》；《后汉书·班超传》；《后汉纪》卷五十八。这些材料较适合钉住事件的发生、官方表述或特定人物、地点，却不自动构成对全境人口、收入、兵数、信仰或动机的调查。账本的置信与材料提示是“遣使事实可靠；到达地点、海路信息和“大秦”识别有争议”；来源限度进一步说明“甘英西使、安息中介与欧亚交通信息研究”。因此，不能把条目中的政权名称、行政分类或战报修辞直接转译为固定族群和统一社会意志，也不能将制度宣布写成各地同步实施。

在叙事上，更稳妥的做法是把这一事件视为一条需要地方中介才能运作的链：中央的命令、地方官与精英的选择、运输与劳力的条件、以及普通家庭可能面对的迁徙、赋役或安全风险彼此相连。现有材料能够显示其中若干环节，却保留了大量沉默。承认这种沉默并非放弃解释，而是避免以一条编年记录覆盖不同区域和社群的差异。

\eventanchor{EH-0102-BAN-CHAO-DEATH}{永元十四年（102年）·班超返洛阳后去世，东汉西域经营失去长期协调者}账本条目\texttt{EH-0102-BAN-CHAO-DEATH}记载，永元十四年（102年），东汉与西域诸国发生“班超返洛阳后去世，东汉西域经营失去长期协调者”。这一节点可放在班超长期依靠地方盟友、军队和个人威望维系都护关系，年老请求返京的条件下理解；其后续影响被账本概括为西域体系的脆弱性更加显现，后任难以完全接续其个人网络，东汉投入与收益持续受质疑。就地方尺度而言，事件并不只属于朝廷或统帅：征发、道路、仓储、户籍、市场、寺院与家庭关系都可能决定命令如何抵达、战事如何延长，或接收何以在不同地区呈现不一样的速度。

本条列举的核心材料为《后汉书·班超传》；《后汉书·和帝纪》；《后汉纪》卷六十三。这些材料较适合钉住事件的发生、官方表述或特定人物、地点，却不自动构成对全境人口、收入、兵数、信仰或动机的调查。账本的置信与材料提示是“归还与卒年可靠；其离任对各绿洲政治的直接影响有争议”；来源限度进一步说明“班超身后西域、个人网络与都护制度研究”。因此，不能把条目中的政权名称、行政分类或战报修辞直接转译为固定族群和统一社会意志，也不能将制度宣布写成各地同步实施。

在叙事上，更稳妥的做法是把这一事件视为一条需要地方中介才能运作的链：中央的命令、地方官与精英的选择、运输与劳力的条件、以及普通家庭可能面对的迁徙、赋役或安全风险彼此相连。现有材料能够显示其中若干环节，却保留了大量沉默。承认这种沉默并非放弃解释，而是避免以一条编年记录覆盖不同区域和社群的差异。

\eventanchor{EH-0105-HEDI-DEATH}{元兴元年（105年）·和帝去世，襁褓中的刘隆即位，是为殇帝，邓太后临朝}账本条目\texttt{EH-0105-HEDI-DEATH}记载，元兴元年（105年），东汉发生“和帝去世，襁褓中的刘隆即位，是为殇帝，邓太后临朝”。这一节点可放在和帝早逝且继承人年幼，邓氏以皇太后身份掌握诏令、禁军和人事渠道的条件下理解；其后续影响被账本概括为东汉再次进入外戚摄政，皇帝年龄使宫廷派系和地方危机更难由常规君主权威调节。就地方尺度而言，事件并不只属于朝廷或统帅：征发、道路、仓储、户籍、市场、寺院与家庭关系都可能决定命令如何抵达、战事如何延长，或接收何以在不同地区呈现不一样的速度。

本条列举的核心材料为《后汉书·和帝纪》；《后汉书·殇帝纪》；《后汉书·邓皇后纪》。这些材料较适合钉住事件的发生、官方表述或特定人物、地点，却不自动构成对全境人口、收入、兵数、信仰或动机的调查。账本的置信与材料提示是“继位年月可靠；和帝病逝与继嗣选择的宫廷过程有争议”；来源限度进一步说明“和帝身后继承、婴儿皇帝与邓氏摄政研究”。因此，不能把条目中的政权名称、行政分类或战报修辞直接转译为固定族群和统一社会意志，也不能将制度宣布写成各地同步实施。

在叙事上，更稳妥的做法是把这一事件视为一条需要地方中介才能运作的链：中央的命令、地方官与精英的选择、运输与劳力的条件、以及普通家庭可能面对的迁徙、赋役或安全风险彼此相连。现有材料能够显示其中若干环节，却保留了大量沉默。承认这种沉默并非放弃解释，而是避免以一条编年记录覆盖不同区域和社群的差异。

\eventanchor{EH-0106-AN-DI-ACCESSION}{延平元年（106年）·殇帝去世，清河王刘祜即位，是为安帝，邓太后继续临朝}账本条目\texttt{EH-0106-AN-DI-ACCESSION}记载，延平元年（106年），东汉发生“殇帝去世，清河王刘祜即位，是为安帝，邓太后继续临朝”。这一节点可放在婴儿殇帝夭亡，朝廷需在刘氏宗室中选择可被邓氏与大臣接受的继承人的条件下理解；其后续影响被账本概括为安帝即位维持皇统连续，却使邓氏摄政延长，外戚权力与皇帝成年后的关系埋下冲突。就地方尺度而言，事件并不只属于朝廷或统帅：征发、道路、仓储、户籍、市场、寺院与家庭关系都可能决定命令如何抵达、战事如何延长，或接收何以在不同地区呈现不一样的速度。

本条列举的核心材料为《后汉书·殇帝纪》；《后汉书·安帝纪》；《后汉书·邓皇后纪》。这些材料较适合钉住事件的发生、官方表述或特定人物、地点，却不自动构成对全境人口、收入、兵数、信仰或动机的调查。账本的置信与材料提示是“继统事实可靠；废立候选与邓氏决策过程有争议”；来源限度进一步说明“安帝继统、邓太后摄政与宗室选择研究”。因此，不能把条目中的政权名称、行政分类或战报修辞直接转译为固定族群和统一社会意志，也不能将制度宣布写成各地同步实施。

在叙事上，更稳妥的做法是把这一事件视为一条需要地方中介才能运作的链：中央的命令、地方官与精英的选择、运输与劳力的条件、以及普通家庭可能面对的迁徙、赋役或安全风险彼此相连。现有材料能够显示其中若干环节，却保留了大量沉默。承认这种沉默并非放弃解释，而是避免以一条编年记录覆盖不同区域和社群的差异。

\eventanchor{EH-0107-LIANG-QIANG-WAR}{永初元年（107年）·西羌在凉州、并州一带再起，东汉长期投入边郡军费与移民安置}账本条目\texttt{EH-0107-LIANG-QIANG-WAR}记载，永初元年（107年），东汉与西羌发生“西羌在凉州、并州一带再起，东汉长期投入边郡军费与移民安置”。这一节点可放在官府征发、边郡冲突和羌部联盟变化使此前平叛未能消除结构性矛盾的条件下理解；其后续影响被账本概括为凉州成为持续军事财政负担，流民、降附与迁徙改变地方社会，正史军功数字须谨慎使用。就地方尺度而言，事件并不只属于朝廷或统帅：征发、道路、仓储、户籍、市场、寺院与家庭关系都可能决定命令如何抵达、战事如何延长，或接收何以在不同地区呈现不一样的速度。

本条列举的核心材料为《后汉书·安帝纪》；《后汉书·西羌传》；《后汉纪》卷六十八。这些材料较适合钉住事件的发生、官方表述或特定人物、地点，却不自动构成对全境人口、收入、兵数、信仰或动机的调查。账本的置信与材料提示是“战乱开始年份可靠；参战部落和各郡受害程度有争议”；来源限度进一步说明“永初羌乱、凉州财政与边民迁徙研究”。因此，不能把条目中的政权名称、行政分类或战报修辞直接转译为固定族群和统一社会意志，也不能将制度宣布写成各地同步实施。

在叙事上，更稳妥的做法是把这一事件视为一条需要地方中介才能运作的链：中央的命令、地方官与精英的选择、运输与劳力的条件、以及普通家庭可能面对的迁徙、赋役或安全风险彼此相连。现有材料能够显示其中若干环节，却保留了大量沉默。承认这种沉默并非放弃解释，而是避免以一条编年记录覆盖不同区域和社群的差异。

\eventanchor{EH-0116-DENG-SUI-DEATH}{元初三年（116年）·邓太后去世，安帝开始亲政并贬逐部分邓氏成员}账本条目\texttt{EH-0116-DENG-SUI-DEATH}记载，元初三年（116年），东汉发生“邓太后去世，安帝开始亲政并贬逐部分邓氏成员”。这一节点可放在安帝成年后长期受邓氏限制，太后去世移除了维持摄政平衡的核心人物的条件下理解；其后续影响被账本概括为皇帝重新控制人事，邓氏网络被削弱；外戚摄政的退出再次以报复性重组而非制度化交接收场。就地方尺度而言，事件并不只属于朝廷或统帅：征发、道路、仓储、户籍、市场、寺院与家庭关系都可能决定命令如何抵达、战事如何延长，或接收何以在不同地区呈现不一样的速度。

本条列举的核心材料为《后汉书·安帝纪》；《后汉书·邓皇后纪》；《后汉纪》卷七十七。这些材料较适合钉住事件的发生、官方表述或特定人物、地点，却不自动构成对全境人口、收入、兵数、信仰或动机的调查。账本的置信与材料提示是“太后卒年与政局转换可靠；安帝报复的范围和个案罪名有争议”；来源限度进一步说明“邓太后身后、安帝亲政与外戚清算研究”。因此，不能把条目中的政权名称、行政分类或战报修辞直接转译为固定族群和统一社会意志，也不能将制度宣布写成各地同步实施。

在叙事上，更稳妥的做法是把这一事件视为一条需要地方中介才能运作的链：中央的命令、地方官与精英的选择、运输与劳力的条件、以及普通家庭可能面对的迁徙、赋役或安全风险彼此相连。现有材料能够显示其中若干环节，却保留了大量沉默。承认这种沉默并非放弃解释，而是避免以一条编年记录覆盖不同区域和社群的差异。

\eventanchor{EH-0125-AN-DI-SUCCESSION}{延光四年（125年）·安帝去世，阎后先立北乡侯，宦官孙程等旋迎立济阴王刘保为顺帝}账本条目\texttt{EH-0125-AN-DI-SUCCESSION}记载，延光四年（125年），东汉发生“安帝去世，阎后先立北乡侯，宦官孙程等旋迎立济阴王刘保为顺帝”。这一节点可放在安帝遗嗣与阎后集团的利益冲突使继承程序分裂，阎氏缺少足以压制宫中近侍的联盟的条件下理解；其后续影响被账本概括为顺帝即位，宦官以拥立功获得封赏和政治资本，宫廷继承更依赖禁中武力与信息通道。就地方尺度而言，事件并不只属于朝廷或统帅：征发、道路、仓储、户籍、市场、寺院与家庭关系都可能决定命令如何抵达、战事如何延长，或接收何以在不同地区呈现不一样的速度。

本条列举的核心材料为《后汉书·安帝纪》；《后汉书·顺帝纪》；《后汉书·孙程传》。这些材料较适合钉住事件的发生、官方表述或特定人物、地点，却不自动构成对全境人口、收入、兵数、信仰或动机的调查。账本的置信与材料提示是“继位转折可靠；北乡侯死因和孙程等人行动细节有争议”；来源限度进一步说明“延光政变、阎氏与宦官拥立研究”。因此，不能把条目中的政权名称、行政分类或战报修辞直接转译为固定族群和统一社会意志，也不能将制度宣布写成各地同步实施。

在叙事上，更稳妥的做法是把这一事件视为一条需要地方中介才能运作的链：中央的命令、地方官与精英的选择、运输与劳力的条件、以及普通家庭可能面对的迁徙、赋役或安全风险彼此相连。现有材料能够显示其中若干环节，却保留了大量沉默。承认这种沉默并非放弃解释，而是避免以一条编年记录覆盖不同区域和社群的差异。

\eventanchor{EH-0144-SHUNDI-DEATH}{建康元年（144年）·顺帝去世，幼帝刘炳即位，是为冲帝，梁太后与梁冀临朝}账本条目\texttt{EH-0144-SHUNDI-DEATH}记载，建康元年（144年），东汉发生“顺帝去世，幼帝刘炳即位，是为冲帝，梁太后与梁冀临朝”。这一节点可放在顺帝留下年幼继承人，梁氏凭皇后、外戚和朝廷官职控制继承安排的条件下理解；其后续影响被账本概括为梁太后和梁冀进入权力中心，外戚对禁军和任官的影响显著扩大。就地方尺度而言，事件并不只属于朝廷或统帅：征发、道路、仓储、户籍、市场、寺院与家庭关系都可能决定命令如何抵达、战事如何延长，或接收何以在不同地区呈现不一样的速度。

本条列举的核心材料为《后汉书·顺帝纪》；《后汉书·冲帝纪》；《后汉书·梁冀传》。这些材料较适合钉住事件的发生、官方表述或特定人物、地点，却不自动构成对全境人口、收入、兵数、信仰或动机的调查。账本的置信与材料提示是“继位年月可靠；梁氏在军政中的分工有争议”；来源限度进一步说明“顺帝身后、冲帝与梁氏摄政研究”。因此，不能把条目中的政权名称、行政分类或战报修辞直接转译为固定族群和统一社会意志，也不能将制度宣布写成各地同步实施。

在叙事上，更稳妥的做法是把这一事件视为一条需要地方中介才能运作的链：中央的命令、地方官与精英的选择、运输与劳力的条件、以及普通家庭可能面对的迁徙、赋役或安全风险彼此相连。现有材料能够显示其中若干环节，却保留了大量沉默。承认这种沉默并非放弃解释，而是避免以一条编年记录覆盖不同区域和社群的差异。

\eventanchor{EH-0145-CHONGDI-DEATH}{永嘉元年（145年）·冲帝夭亡，梁太后立八岁的刘缵为质帝，梁冀继续掌政}账本条目\texttt{EH-0145-CHONGDI-DEATH}记载，永嘉元年（145年），东汉发生“冲帝夭亡，梁太后立八岁的刘缵为质帝，梁冀继续掌政”。这一节点可放在冲帝在位短暂，梁氏必须迅速另择宗室幼主以保持摄政结构的条件下理解；其后续影响被账本概括为连续的幼帝继承使皇帝权威进一步空洞化，梁冀的实际权力和朝臣依附关系更加强化。就地方尺度而言，事件并不只属于朝廷或统帅：征发、道路、仓储、户籍、市场、寺院与家庭关系都可能决定命令如何抵达、战事如何延长，或接收何以在不同地区呈现不一样的速度。

本条列举的核心材料为《后汉书·冲帝纪》；《后汉书·质帝纪》；《后汉书·梁冀传》。这些材料较适合钉住事件的发生、官方表述或特定人物、地点，却不自动构成对全境人口、收入、兵数、信仰或动机的调查。账本的置信与材料提示是“冲帝卒与质帝即位可靠；选立动机与宫廷谈判有争议”；来源限度进一步说明“幼帝更替、梁氏摄政与继承脆弱性研究”。因此，不能把条目中的政权名称、行政分类或战报修辞直接转译为固定族群和统一社会意志，也不能将制度宣布写成各地同步实施。

在叙事上，更稳妥的做法是把这一事件视为一条需要地方中介才能运作的链：中央的命令、地方官与精英的选择、运输与劳力的条件、以及普通家庭可能面对的迁徙、赋役或安全风险彼此相连。现有材料能够显示其中若干环节，却保留了大量沉默。承认这种沉默并非放弃解释，而是避免以一条编年记录覆盖不同区域和社群的差异。

\eventanchor{EH-0146-ZHIDI-DEATH-HUANDI}{本初元年（146年）·质帝被毒杀后，梁太后立刘志为桓帝，梁冀继续控制朝政}账本条目\texttt{EH-0146-ZHIDI-DEATH-HUANDI}记载，本初元年（146年），东汉发生“质帝被毒杀后，梁太后立刘志为桓帝，梁冀继续控制朝政”。这一节点可放在质帝对梁冀的公开批评威胁摄政集团，梁氏需寻找较易控制的成年宗室的条件下理解；其后续影响被账本概括为桓帝登位却受梁冀掣肘，毒杀指控成为理解外戚专权的重要叙事，亦需区分史家谴责与可证事实。就地方尺度而言，事件并不只属于朝廷或统帅：征发、道路、仓储、户籍、市场、寺院与家庭关系都可能决定命令如何抵达、战事如何延长，或接收何以在不同地区呈现不一样的速度。

本条列举的核心材料为《后汉书·质帝纪》；《后汉书·桓帝纪》；《后汉书·梁冀传》。这些材料较适合钉住事件的发生、官方表述或特定人物、地点，却不自动构成对全境人口、收入、兵数、信仰或动机的调查。账本的置信与材料提示是“废立与桓帝即位可靠；毒杀责任主要见于后出史家叙事，证据有争议”；来源限度进一步说明“质帝之死、梁冀专权与后汉史叙事批判研究”。因此，不能把条目中的政权名称、行政分类或战报修辞直接转译为固定族群和统一社会意志，也不能将制度宣布写成各地同步实施。

在叙事上，更稳妥的做法是把这一事件视为一条需要地方中介才能运作的链：中央的命令、地方官与精英的选择、运输与劳力的条件、以及普通家庭可能面对的迁徙、赋役或安全风险彼此相连。现有材料能够显示其中若干环节，却保留了大量沉默。承认这种沉默并非放弃解释，而是避免以一条编年记录覆盖不同区域和社群的差异。

\subsection{信仰、知识与物质空间}

\eventanchor{EH-0147-LI-GU-PURGE}{建和元年（147年）·梁冀诬杀李固、杜乔等反对者，朝廷清流官僚受挫}账本条目\texttt{EH-0147-LI-GU-PURGE}记载，建和元年（147年），东汉发生“梁冀诬杀李固、杜乔等反对者，朝廷清流官僚受挫”。这一节点可放在李固等主张限制梁氏并批评连续幼帝政治，触及梁冀对人事和继统的控制的条件下理解；其后续影响被账本概括为反梁官僚被清除，士人以名节评价政治的语言更趋尖锐，朝廷监督机制进一步失衡。就地方尺度而言，事件并不只属于朝廷或统帅：征发、道路、仓储、户籍、市场、寺院与家庭关系都可能决定命令如何抵达、战事如何延长，或接收何以在不同地区呈现不一样的速度。

本条列举的核心材料为《后汉书·李固传》；《后汉书·杜乔传》；《后汉纪》卷八十三。这些材料较适合钉住事件的发生、官方表述或特定人物、地点，却不自动构成对全境人口、收入、兵数、信仰或动机的调查。账本的置信与材料提示是“处置事实可靠；党派标签和密谋情节有争议”；来源限度进一步说明“李固案、外戚专权与官僚反抗研究”。因此，不能把条目中的政权名称、行政分类或战报修辞直接转译为固定族群和统一社会意志，也不能将制度宣布写成各地同步实施。

在叙事上，更稳妥的做法是把这一事件视为一条需要地方中介才能运作的链：中央的命令、地方官与精英的选择、运输与劳力的条件、以及普通家庭可能面对的迁徙、赋役或安全风险彼此相连。现有材料能够显示其中若干环节，却保留了大量沉默。承认这种沉默并非放弃解释，而是避免以一条编年记录覆盖不同区域和社群的差异。

\eventanchor{EH-0159-LIANG-JI-FALL}{延熹二年（159年）·桓帝与单超等宦官诛灭梁冀集团，梁氏外戚专权结束}账本条目\texttt{EH-0159-LIANG-JI-FALL}记载，延熹二年（159年），东汉发生“桓帝与单超等宦官诛灭梁冀集团，梁氏外戚专权结束”。这一节点可放在梁冀长期控制朝政并压制皇帝，桓帝只能借禁中宦官发动突袭的条件下理解；其后续影响被账本概括为梁氏被清算，宦官因功受封掌权；外戚问题未被制度解决，而转化为宦官—士人冲突。就地方尺度而言，事件并不只属于朝廷或统帅：征发、道路、仓储、户籍、市场、寺院与家庭关系都可能决定命令如何抵达、战事如何延长，或接收何以在不同地区呈现不一样的速度。

本条列举的核心材料为《后汉书·桓帝纪》；《后汉书·梁冀传》；《后汉书·宦者列传》。这些材料较适合钉住事件的发生、官方表述或特定人物、地点，却不自动构成对全境人口、收入、兵数、信仰或动机的调查。账本的置信与材料提示是“政变年月可靠；桓帝主导程度与查抄财产数额有争议”；来源限度进一步说明“梁冀覆灭、宦官政变与外戚财产研究”。因此，不能把条目中的政权名称、行政分类或战报修辞直接转译为固定族群和统一社会意志，也不能将制度宣布写成各地同步实施。

在叙事上，更稳妥的做法是把这一事件视为一条需要地方中介才能运作的链：中央的命令、地方官与精英的选择、运输与劳力的条件、以及普通家庭可能面对的迁徙、赋役或安全风险彼此相连。现有材料能够显示其中若干环节，却保留了大量沉默。承认这种沉默并非放弃解释，而是避免以一条编年记录覆盖不同区域和社群的差异。

\eventanchor{EH-0166-FIRST-PARTISAN-PROHIBITION}{延熹九年（166年）·桓帝朝以李膺等结党为名逮捕太学生和官僚，形成第一次党锢}账本条目\texttt{EH-0166-FIRST-PARTISAN-PROHIBITION}记载，延熹九年（166年），东汉发生“桓帝朝以李膺等结党为名逮捕太学生和官僚，形成第一次党锢”。这一节点可放在宦官与依附官僚担忧太学生、名士的舆论和荐举网络挑战其任官利益的条件下理解；其后续影响被账本概括为清议士人与宦官集团的对立公开化；“党人”是国家刑名分类，不能等同于统一政治组织。就地方尺度而言，事件并不只属于朝廷或统帅：征发、道路、仓储、户籍、市场、寺院与家庭关系都可能决定命令如何抵达、战事如何延长，或接收何以在不同地区呈现不一样的速度。

本条列举的核心材料为《后汉书·桓帝纪》；《后汉书·党锢列传》；《后汉纪》卷九十二。这些材料较适合钉住事件的发生、官方表述或特定人物、地点，却不自动构成对全境人口、收入、兵数、信仰或动机的调查。账本的置信与材料提示是“案件与禁锢存在可靠；被捕人数、党人边界和罪名有争议”；来源限度进一步说明“第一次党锢、太学生网络与宦官政治研究”。因此，不能把条目中的政权名称、行政分类或战报修辞直接转译为固定族群和统一社会意志，也不能将制度宣布写成各地同步实施。

在叙事上，更稳妥的做法是把这一事件视为一条需要地方中介才能运作的链：中央的命令、地方官与精英的选择、运输与劳力的条件、以及普通家庭可能面对的迁徙、赋役或安全风险彼此相连。现有材料能够显示其中若干环节，却保留了大量沉默。承认这种沉默并非放弃解释，而是避免以一条编年记录覆盖不同区域和社群的差异。

\eventanchor{EH-0167-DUAN-JIONG-QIANG}{永康元年（167年）至建宁二年（169年）·段颎在陇右连续进军，东羌大部被击散或降附，战事伴随严酷杀戮}账本条目\texttt{EH-0167-DUAN-JIONG-QIANG}记载，永康元年（167年）至建宁二年（169年），东汉与东羌发生“段颎在陇右连续进军，东羌大部被击散或降附，战事伴随严酷杀戮”。这一节点可放在长期羌乱消耗财政并威胁陇右交通，朝廷授权段颎实施高强度围剿的条件下理解；其后续影响被账本概括为短期内恢复部分通道与郡县控制，却加重地方仇怨和军政依赖，后续凉州危机并未根除。就地方尺度而言，事件并不只属于朝廷或统帅：征发、道路、仓储、户籍、市场、寺院与家庭关系都可能决定命令如何抵达、战事如何延长，或接收何以在不同地区呈现不一样的速度。

本条列举的核心材料为《后汉书·段颎传》；《后汉书·西羌传》；《后汉纪》卷九十三至九十五。这些材料较适合钉住事件的发生、官方表述或特定人物、地点，却不自动构成对全境人口、收入、兵数、信仰或动机的调查。账本的置信与材料提示是“战役序列大体可靠；斩获、迁徙和“平定”范围有争议”；来源限度进一步说明“段颎征羌、凉州军事化与战争暴力研究”。因此，不能把条目中的政权名称、行政分类或战报修辞直接转译为固定族群和统一社会意志，也不能将制度宣布写成各地同步实施。

在叙事上，更稳妥的做法是把这一事件视为一条需要地方中介才能运作的链：中央的命令、地方官与精英的选择、运输与劳力的条件、以及普通家庭可能面对的迁徙、赋役或安全风险彼此相连。现有材料能够显示其中若干环节，却保留了大量沉默。承认这种沉默并非放弃解释，而是避免以一条编年记录覆盖不同区域和社群的差异。

\eventanchor{EH-0168-LINGDI-ACCESSION}{永康元年（168年）·桓帝去世，十二岁的刘宏即位，是为灵帝，窦太后、窦武与陈蕃辅政}账本条目\texttt{EH-0168-LINGDI-ACCESSION}记载，永康元年（168年），东汉发生“桓帝去世，十二岁的刘宏即位，是为灵帝，窦太后、窦武与陈蕃辅政”。这一节点可放在桓帝无子且宦官势力已深植禁中，朝廷选择河间王刘开一支的刘宏维持刘氏皇统的条件下理解；其后续影响被账本概括为幼帝即位使外戚、士人大臣与宦官围绕辅政权和禁军控制展开直接竞争。就地方尺度而言，事件并不只属于朝廷或统帅：征发、道路、仓储、户籍、市场、寺院与家庭关系都可能决定命令如何抵达、战事如何延长，或接收何以在不同地区呈现不一样的速度。

本条列举的核心材料为《后汉书·桓帝纪》；《后汉书·灵帝纪》；《后汉书·窦武传》。这些材料较适合钉住事件的发生、官方表述或特定人物、地点，却不自动构成对全境人口、收入、兵数、信仰或动机的调查。账本的置信与材料提示是“继位年月可靠；窦武、陈蕃的职权边界有争议”；来源限度进一步说明“灵帝继统、窦武陈蕃辅政与宫廷权力研究”。因此，不能把条目中的政权名称、行政分类或战报修辞直接转译为固定族群和统一社会意志，也不能将制度宣布写成各地同步实施。

在叙事上，更稳妥的做法是把这一事件视为一条需要地方中介才能运作的链：中央的命令、地方官与精英的选择、运输与劳力的条件、以及普通家庭可能面对的迁徙、赋役或安全风险彼此相连。现有材料能够显示其中若干环节，却保留了大量沉默。承认这种沉默并非放弃解释，而是避免以一条编年记录覆盖不同区域和社群的差异。

\eventanchor{EH-0168-DOU-WU-COUP}{建宁元年（168年）·宦官曹节、王甫等发动政变，杀窦武、陈蕃，窦太后被幽禁}账本条目\texttt{EH-0168-DOU-WU-COUP}记载，建宁元年（168年），东汉发生“宦官曹节、王甫等发动政变，杀窦武、陈蕃，窦太后被幽禁”。这一节点可放在窦武、陈蕃谋诛宦官但未能稳控宫门、诏令和北军，宦官利用皇帝近侍地位反击的条件下理解；其后续影响被账本概括为外戚—士人联合失败，宦官取得压倒性禁中优势，灵帝朝政治空间急剧收缩。就地方尺度而言，事件并不只属于朝廷或统帅：征发、道路、仓储、户籍、市场、寺院与家庭关系都可能决定命令如何抵达、战事如何延长，或接收何以在不同地区呈现不一样的速度。

本条列举的核心材料为《后汉书·灵帝纪》；《后汉书·窦武传》；《后汉书·陈蕃传》。这些材料较适合钉住事件的发生、官方表述或特定人物、地点，却不自动构成对全境人口、收入、兵数、信仰或动机的调查。账本的置信与材料提示是“政变事实可靠；先发制人责任与战斗细节有争议”；来源限度进一步说明“窦武陈蕃败亡、禁军与宦官政变研究”。因此，不能把条目中的政权名称、行政分类或战报修辞直接转译为固定族群和统一社会意志，也不能将制度宣布写成各地同步实施。

在叙事上，更稳妥的做法是把这一事件视为一条需要地方中介才能运作的链：中央的命令、地方官与精英的选择、运输与劳力的条件、以及普通家庭可能面对的迁徙、赋役或安全风险彼此相连。现有材料能够显示其中若干环节，却保留了大量沉默。承认这种沉默并非放弃解释，而是避免以一条编年记录覆盖不同区域和社群的差异。

\eventanchor{EH-0169-SECOND-PARTISAN-PROHIBITION}{建宁二年（169年）·宦官集团扩大逮捕、禁锢党人，第二次党锢形成}账本条目\texttt{EH-0169-SECOND-PARTISAN-PROHIBITION}记载，建宁二年（169年），东汉发生“宦官集团扩大逮捕、禁锢党人，第二次党锢形成”。这一节点可放在窦武政变后宦官将士人批评视为潜在复辟威胁，以结党罪名清除官学和地方声望网络的条件下理解；其后续影响被账本概括为大量士人失去仕进资格，中央与地方精英沟通受损，后期危机中清议人物将以不同方式重返政治。就地方尺度而言，事件并不只属于朝廷或统帅：征发、道路、仓储、户籍、市场、寺院与家庭关系都可能决定命令如何抵达、战事如何延长，或接收何以在不同地区呈现不一样的速度。

本条列举的核心材料为《后汉书·灵帝纪》；《后汉书·党锢列传》；《资治通鉴》卷五十六。这些材料较适合钉住事件的发生、官方表述或特定人物、地点，却不自动构成对全境人口、收入、兵数、信仰或动机的调查。账本的置信与材料提示是“禁锢发生可靠；牵连名单、地域规模和持续期限有争议”；来源限度进一步说明“第二次党锢、清议网络与政治排斥研究”。因此，不能把条目中的政权名称、行政分类或战报修辞直接转译为固定族群和统一社会意志，也不能将制度宣布写成各地同步实施。

在叙事上，更稳妥的做法是把这一事件视为一条需要地方中介才能运作的链：中央的命令、地方官与精英的选择、运输与劳力的条件、以及普通家庭可能面对的迁徙、赋役或安全风险彼此相连。现有材料能够显示其中若干环节，却保留了大量沉默。承认这种沉默并非放弃解释，而是避免以一条编年记录覆盖不同区域和社群的差异。

\eventanchor{EH-0175-XIPING-STONE-CLASSICS}{熹平四年（175年）·蔡邕等校定儒家经典并刻置熹平石经于洛阳太学}账本条目\texttt{EH-0175-XIPING-STONE-CLASSICS}记载，熹平四年（175年），东汉发生“蔡邕等校定儒家经典并刻置熹平石经于洛阳太学”。这一节点可放在经文异字和师法竞争影响教学与官学权威，朝廷希望用可公开校读的石刻稳定文本的条件下理解；其后续影响被账本概括为石经为太学与经学权威提供物质锚点；残石和著录可校字形，不能完整复原所有教学实践。就地方尺度而言，事件并不只属于朝廷或统帅：征发、道路、仓储、户籍、市场、寺院与家庭关系都可能决定命令如何抵达、战事如何延长，或接收何以在不同地区呈现不一样的速度。

本条列举的核心材料为《后汉书·蔡邕传》；《后汉书·灵帝纪》；熹平石经残石及洛阳太学相关考古材料。这些材料较适合钉住事件的发生、官方表述或特定人物、地点，却不自动构成对全境人口、收入、兵数、信仰或动机的调查。账本的置信与材料提示是“立石经事实可靠；所刻篇目、完成阶段和现存残石归属有争议”；来源限度进一步说明“熹平石经、经学定本与太学物质文化研究”。因此，不能把条目中的政权名称、行政分类或战报修辞直接转译为固定族群和统一社会意志，也不能将制度宣布写成各地同步实施。

在叙事上，更稳妥的做法是把这一事件视为一条需要地方中介才能运作的链：中央的命令、地方官与精英的选择、运输与劳力的条件、以及普通家庭可能面对的迁徙、赋役或安全风险彼此相连。现有材料能够显示其中若干环节，却保留了大量沉默。承认这种沉默并非放弃解释，而是避免以一条编年记录覆盖不同区域和社群的差异。

\eventanchor{EH-0177-XIANBEI-RAIDS}{熹平六年（177年）·檀石槐领导下的鲜卑多次侵扰北边，汉军出击受挫}账本条目\texttt{EH-0177-XIANBEI-RAIDS}记载，熹平六年（177年），东汉与鲜卑诸部发生“檀石槐领导下的鲜卑多次侵扰北边，汉军出击受挫”。这一节点可放在北匈奴衰退后鲜卑部众扩大活动空间，东汉边郡在羌乱和财政紧张中难以维持机动防御的条件下理解；其后续影响被账本概括为北方威胁由匈奴转向多部鲜卑联盟，边郡军费与募兵压力进一步侵蚀中央财政。就地方尺度而言，事件并不只属于朝廷或统帅：征发、道路、仓储、户籍、市场、寺院与家庭关系都可能决定命令如何抵达、战事如何延长，或接收何以在不同地区呈现不一样的速度。

本条列举的核心材料为《后汉书·鲜卑传》；《后汉书·灵帝纪》；《后汉纪》卷一百零三。这些材料较适合钉住事件的发生、官方表述或特定人物、地点，却不自动构成对全境人口、收入、兵数、信仰或动机的调查。账本的置信与材料提示是“边患与出师事实可靠；鲜卑政治统一程度、兵数和损失有争议”；来源限度进一步说明“檀石槐、鲜卑联盟与东汉北边防务研究”。因此，不能把条目中的政权名称、行政分类或战报修辞直接转译为固定族群和统一社会意志，也不能将制度宣布写成各地同步实施。

在叙事上，更稳妥的做法是把这一事件视为一条需要地方中介才能运作的链：中央的命令、地方官与精英的选择、运输与劳力的条件、以及普通家庭可能面对的迁徙、赋役或安全风险彼此相连。现有材料能够显示其中若干环节，却保留了大量沉默。承认这种沉默并非放弃解释，而是避免以一条编年记录覆盖不同区域和社群的差异。

\eventanchor{EH-0178-HONGDU-GATE}{光和元年（178年）·灵帝置鸿都门学，招集辞赋、书画与尺牍之士入学}账本条目\texttt{EH-0178-HONGDU-GATE}记载，光和元年（178年），东汉发生“灵帝置鸿都门学，招集辞赋、书画与尺牍之士入学”。这一节点可放在宦官和皇帝试图绕开太学清议与经师网络，培植较少依赖旧式荐举的近侍文化资源的条件下理解；其后续影响被账本概括为文学与书画技能获得官方位置，也加剧士人对卖官、近幸和选官标准的批评。就地方尺度而言，事件并不只属于朝廷或统帅：征发、道路、仓储、户籍、市场、寺院与家庭关系都可能决定命令如何抵达、战事如何延长，或接收何以在不同地区呈现不一样的速度。

本条列举的核心材料为《后汉书·灵帝纪》；《后汉书·蔡邕传》；《后汉纪》卷一百零四。这些材料较适合钉住事件的发生、官方表述或特定人物、地点，却不自动构成对全境人口、收入、兵数、信仰或动机的调查。账本的置信与材料提示是“设置事实可靠；学员身份、课程与政治功能有争议”；来源限度进一步说明“鸿都门学、文学技艺与灵帝朝用人研究”。因此，不能把条目中的政权名称、行政分类或战报修辞直接转译为固定族群和统一社会意志，也不能将制度宣布写成各地同步实施。

在叙事上，更稳妥的做法是把这一事件视为一条需要地方中介才能运作的链：中央的命令、地方官与精英的选择、运输与劳力的条件、以及普通家庭可能面对的迁徙、赋役或安全风险彼此相连。现有材料能够显示其中若干环节，却保留了大量沉默。承认这种沉默并非放弃解释，而是避免以一条编年记录覆盖不同区域和社群的差异。

\eventanchor{EH-0183-TAIPING-DAO}{光和六年（183年）前后·张角以太平道组织信众、设置方主并传播“苍天已死”动员口号}账本条目\texttt{EH-0183-TAIPING-DAO}记载，光和六年（183年）前后，冀州太平道集团发生“张角以太平道组织信众、设置方主并传播“苍天已死”动员口号”。这一节点可放在灾害、赋役、地方治理失效与被禁锢士人无法入仕共同扩大了非正式结社的空间的条件下理解；其后续影响被账本概括为跨郡信众网络将宗教治疗和政治预言结合，为次年大规模起事提供动员基础；口号文本有后世定型可能。就地方尺度而言，事件并不只属于朝廷或统帅：征发、道路、仓储、户籍、市场、寺院与家庭关系都可能决定命令如何抵达、战事如何延长，或接收何以在不同地区呈现不一样的速度。

本条列举的核心材料为《后汉书·皇甫嵩传》；《后汉书·张角传》；《后汉纪》卷一百零九。这些材料较适合钉住事件的发生、官方表述或特定人物、地点，却不自动构成对全境人口、收入、兵数、信仰或动机的调查。账本的置信与材料提示是“组织活动大体可靠；成员数量、教义文本和起事准备程度有争议”；来源限度进一步说明“太平道组织、疾病救治与黄巾动员研究”。因此，不能把条目中的政权名称、行政分类或战报修辞直接转译为固定族群和统一社会意志，也不能将制度宣布写成各地同步实施。

在叙事上，更稳妥的做法是把这一事件视为一条需要地方中介才能运作的链：中央的命令、地方官与精英的选择、运输与劳力的条件、以及普通家庭可能面对的迁徙、赋役或安全风险彼此相连。现有材料能够显示其中若干环节，却保留了大量沉默。承认这种沉默并非放弃解释，而是避免以一条编年记录覆盖不同区域和社群的差异。
